%% ============================================================================
%% V. Methodology                                             target: ~1.3 pages
%% Stage 1: the methodology must be fully developed, including the specific
%% design choices in the system architecture and the experiments.
%% ============================================================================
\section{Methodology}
\label{sec:method}

\begin{figure}[t]
  \centering
  \figplaceholder{\linewidth}{3.2cm}{%
    \textbf{Fig.~3 (pipeline).} Inputs $\gpose$, $\wrench$, $\cpt$
    $\rightarrow$ tip-pose estimator $\rightarrow$ $\tposehat$ $\rightarrow$
    contact-aware steering $\rightarrow$ gripper motion command.}
  \caption{Overview of the approach.}
  \label{fig:pipeline}
\end{figure}

\subsection{Contact as an ad-hoc joint}
\label{sec:method:reach}
\begin{outline}
  \item Model the rod as a quasi-static Kirchhoff rod~\cite{bretl2014quasistatic};
    the edge contact is a point constraint that transmits force but no moment,
    \ie a passive joint whose location along the rod, $s_{\mathrm{c}}$, changes
    as the rod is fed over the edge.
  \item Compute $\Rfree$ and $\Rcon$ by sampling gripper poses in simulation
    (\eg with discrete elastic rods~\cite{bergou2008der}) and validate on
    hardware; this tests \ref{hyp:reach}.
  \item \todo{Analytic result for the planar case (elastica)? Bound on the
    achievable change in tip direction as a function of clearance.}
\end{outline}

\subsection{Sensorized edge}
\label{sec:method:edge}
\begin{outline}
  \item Conductive TPU ring around the opening that reports contact and the
    angular contact location $\cang$ along the ring.
    \todo{Sensing principle, resolution, latency, calibration.}
  \item Design choice: measure the contact location directly instead of
    inferring it from kinematics or wrench (\cf~\cite{bajo2012kinematics}), to
    isolate the value of contact information (\ref{hyp:estimate}).
\end{outline}

\subsection{Haptic tip-pose estimation}
\label{sec:method:estimation}
\begin{outline}
  \item Learn $f_{\theta}:\vect{o}_t\mapsto\tposehat$ from data collected on
    the real system (\cref{sec:experiments}); camera ground truth is used only
    as training labels.
  \item Inputs: $\gpose$, the tared and gravity-compensated wrench $\wrench$,
    the contact state, and $\cang$. Outputs: tip position and direction
    ($(x,y,\theta)$ in the planar case).
  \item Candidate models to compare: MLP; a temporal model over a short history
    (\eg 1D-CNN or GRU) to capture friction and hysteresis; a physics-informed
    model that learns a residual on top of a rod-model
    prediction~\cite{mishani2022realtime}. \todo{Final list of models.}
  \item Loss: weighted position and orientation error; uncertainty from an
    ensemble to gate the controller. \todo{Decide.}
\end{outline}

\subsection{Contact-aware steering}
\label{sec:method:steering}
\begin{outline}
  \item Phases: (i) approach and make contact at a chosen edge location;
    (ii) pivot about the contact to rotate the tip direction; (iii) feed the
    rod to translate the tip to $\tposegoal$.
  \item Closed loop on $\tposehat$: local Jacobian-based servoing (Jacobian
    estimated online or from the learned model), with limits on contact force
    and rod curvature. \todo{Alternative: sampling-based planning in simulation
    with feedback correction.}
\end{outline}

\subsection{Implementation}
\label{sec:method:implementation}
\begin{outline}
  \item ROS~1 (Noetic) with MoveIt for motion planning; synchronized logging of
    joint states, wrench, contact signal, and camera frames.
    \todo{Synchronization method and rates.}
\end{outline}
